\section{Obiettivi}

L’obiettivo è sviluppare un’applicazione \textit{web responsive} in grado di monitorare e di eseguire le azioni sotto menzionate in un sistema di illuminazione pubblico.

\subsection{Assunti preventivi}

\begin{itemize}
    \item il gestore è in possesso di uno smartphone (Android o iOS) e nelle condizioni di utilizzare un browser web.
\end{itemize}

\subsection{Obiettivi obbligatori}

\begin{itemize}
    \item Rilevamento della presenza di persone in prossimità della fonte luminosa: attraverso sensori eterogenei (bluetooth, IR, videocamere, ecc.) il sistema rileva la presenza di persone in un’area della città gestita dal sistema e attiva le luci di conseguenza.
    \item Aumento/riduzione dell’intensità luminosa: il sistema deve essere in grado di comunicare attraverso un protocollo IoT l’aumento o la riduzione dell’intensità luminosa emessa da uno specifico impianto di illuminazione, permettendo una regolazione singola o per l’intera via/piazza (area coperta dal servizio).
    \item Rilevamento automatico del guasto di un impianto di illuminazione: il sistema deve rilevare ogni eventuale malfunzionamento dell’impianto di illuminazione, notificando il gestore ed aprendo un ticket di assistenza su una piattaforma esterna che permetta di identificare la località del guasto.
    \item Segnalazione manuale del guasto di un impianto di illuminazione: il sistema deve permettere di inserire manualmente un nuovo guasto ad un impianto di illuminazione.
    \item Aumento/riduzione manuale dell’intensità luminosa: il sistema deve permette al gestore di aumentare/ridurre o riportare in modalità automatica l’illuminazione di un singolo impianto luminoso o di un’intera area coperta dal servizio.
    \item Inserimento e gestione di un impianto luminoso: il sistema deve permettere al gestore
    l’inserimento di nuovi impianti luminosi e il loro inserimento o rimozione all’interno di un’area coperta dal servizio.
    \item Aumento o riduzione globale dell’intensità luminosa: il sistema deve permettere al gestore l’aumento o la riduzione globale (su tutti gli impianti controllati dell’intensità luminosa, opzione particolarmente utile in condizioni di crepuscolo o luce lunare particolarmente intensa.
    
\end{itemize}

\subsection{Obiettivi secondari accennati con l'azienda}

\begin{itemize}
    \item modulare la luminosità a seconda della quantità di persone o tipologia di entità;
    \item Rilevamento tipologia guasti (Guasti locali o distribuiti, guasti a sistemi di alimentazione);
    \item Preset di illuminazione per le varie situazioni;
    \item integrazione con previsioni di eventi normali(alba e tramonto) ed anomali come eclissi;
    \item Gestione multi utente ("gestore illuminazione pubblica", "installatore/manutentore", "Verificatore d'impianto" , altro);
    \item Pagina disponibile ai "cittadini" dove questi possano (seguendo 2/3 passaggi, ad esempio caricando un'immagine del guasto) caricare una segnalazione di guasto.
\end{itemize}

\subsection{Ottica di espansione futura}

\begin{itemize}
    \item Integrazioni future con sistemi di alimentazione/UPS vari nella gestione guasti/distribuzione dell'alimentazione;
    \item Api pubblica regolamentata utilizzabile in futuro da pannellistica a led per mostrare ai cittadini i risparmi o altro.
\end{itemize}

\section{Prima stesura architetturale}

Visto l'ambito critico di operazione del prodotto software per noi i capostipiti nello sviluppo di questo progetto saranno: alta modularità, alta scalabilità, alta resilienza e facile estensibilità dello stesso.

L'architettura seguirà quindi i principi base di un sistema a microservizi che tenga conto di tutte queste importanti questioni.

\subsection{User stories / casi d'uso}

1. login e logout di un operatore
2. collegamento di un impianto luminoso ai dati derivanti da un sensore (modalità automatica)
3. gestione manuale di un impianto luminoso
4. aumento o riduzione globale dell’intensità luminosa da parte di un operatore o tramite dati di un
sensore
5. aumento o riduzione locale (per area illuminata) dell’intensità luminosa da parte di un operatore o tramite dati di un sensore
6. Gestione ticketing guasti
7. inserimento e gestione di un impianto luminoso
8. creazione, modifica e rimozione di nuove aree illuminate
9. tracciamento delle intensità luminose di ogni impianto.
10. Rilevamento della presenza in un’area illuminata e aumento automatico dell’intensità
luminosa

\subsection{Utenti}

Molteplici saranno gli utenti che utilizzeranno il sistema.


|Utente| utilizzi | Tipo di requisito|
|---|---|---|
|Semplice cittadino|Può vedere una dashboard relativa all'illuminazione| Aggiuntivo|
|Gestore dell'illuminazione| Può impostare l'illuminazione| Obbligatorio|
|Gestore momentaneo| Può impostare l'illuminazione per un periodo limitato di tempo|Aggiuntivo |
|Installatore/manutentore|Aggiunge nuove sezioni illuminanti, risolve i guasti|Obbligatorio|
|Verificatore di impianto|Gira a controllare periodicamente se ci sono guasti ai corpi illuminanti| Obbligatorio|

\subsection{Servizi}

|Servizio|Scopo|Tecnologia|
|---|---|---|
|Mqtt|Comunicazione delle componenti IoT |Mosquitto|
|Database|Stoccaggio a lungo termine dei dati per un'analisi futura degli stessi per prevenire guasti o risolvere situazioni ricorrenti| Postgres|
|Coordinatore|Coordinamento e gestione diretta degli apparati illuminanti| Python|
|ApiREST del sistema d'illuminazione|Api Backend per la webapp e altri utilizzi futuri|Python, Flask|
|Backend/Api Ticketing|Gestisce il sistema di ticketing dei guasti|Python, Flask|
|WebApp|Consente agli utenti(Definiti più avanti) di interfacciarsi con il sistema|VueJs o React|

\subsection{Gestione di Deploy}

Per il deploy sarà utilizzato **Docker** per consentire, alla bisogna, lo scalare orizzontale del sistema, così da poter gestire più utenti abbattendo i costi.
        
Il sistema sarà \textit{multi-tennant} as a service oppure installabile on premise.

\section{Testing}

Ognuno dei servizi avrà la sua specifica strategia di testing.

I test di ognuno dovranno avere comunque almeno l'80% di code coverage e dovranno essere correlati di report relativamente all'esecuzione degli stessi.

\section{Processo di Coordinamento}

\subsection{Scopo}
Di seguito vengono descritte le modalità di coordinamento del gruppo per quanto riguarda la comunicazione, le riunioni, i ruoli di progetto e l'assegnazione dei compiti. Inoltre saranno descritti gli strumenti scelti.

\subsection{Comunicazione}
Durante lo svolgimento del progetto il gruppo \textbf{\textsl{SWEasabi}} comunicherà seguendo le seguenti procedure:

\begin{itemize}
\item \textbf{\textsl{Comunicazioni interne}}
\begin{itemize} 
    \item Le comunicazioni interne avvengono attraverso un canale Discord e un gruppo Telegram. Entrambi scelti per la loro versatilià e per il fatto di essere multi-piattaforma. In particolare, Telegram sarà utilizzato per le comunicazioni più rapide.\end{itemize}

\item \textbf{\textsl{Comunicazioni esterne}}
\begin{itemize} 
    \item Le comunicazioni esterne saranno gestite dal Responsabile di Progetto attraverso un indirizzo di posta elettronica creato appositamente, ovvero \textbf{sweasabi@gmail.com}
\end{itemize}
\end{itemize}

\subsection{Riunioni}
Le riunioni saranno organizzate dal Responsabile di Progetto, il quale nominerà un segretario che avrà il compito di scrivere un \textbf{verbale di riunione}, che dopo essere approvato renderà valida la riunione.
Le riunioni saranno di due tipi:


\item \textbf{Riunioni interne}
\begin{itemize} 
    \item La partecipazione sarà permessa solo ai membri del gruppo, e per essere validata dovrà essere di almeno quattro persone, perché per prendere decisioni è necessaria la maggioranza. La data della riunione verrà fissata dal Responsabile di Progetto, dopo aver verificato la disponibilità dei membri del gruppo, i quali dovranno essere puntuali o avvertire con anticipo in caso di assenze o ritardi.Ogni membro del gruppo potrà richiedere una riunione al Responsabile di Progetto, il quale la organizzerà, se necessario. Per le riunioni interne del gruppo, si è deciso di utilizzare Discord.
\end{itemize}

\item \textbf{Riunioni esterne}
    \begin{itemize} \item La partecipazione sarà permessa al proponente, i committenti e al gruppo *nomegruppo*.
    (descrizione di come/dove, in base al proponente)
    \end{itemize}
\end{itemize}

\subsection{Strumenti}
\begin{center}
    \begin{tabularx}{\textwidth} {
        | >{\raggedright\arraybackslash}>{\hsize=.4\hsize}X |
          >{\raggedright\arraybackslash}>{\hsize=.6\hsize}X |
        }

        \hline
        \textbf{Uso} & \textbf{Strumento} \\
        \hline\hline

        %Uso
        ITS &
        %Strumento
       GitHub  \\
        \hline

        %Uso
        Comunicazioni rapide &
        %Strumento
       Telegram  \\
        \hline

        %Uso
        Comunicazioni visive &
        %Strumento
       Telegram  \\
        \hline

        %Uso
        VCS &
        %Strumento
       GitHub  \\
        \hline

        %Uso
        Presentazioni &
        %Strumento
       Google Slides  \\
        \hline

    \end{tabularx}
\end{center}

\begin{itemize}
\item \textbf{\textsl{ITS}}
\begin{itemize} \item L'ITS ovvero Issue Tracking System viene utilizzato per definire, tracciare e gestire le varie questioni da risolvere.\end{itemize}

\item \textbf{textsl{Comunicazioni rapide}}
\begin{itemize} \item Per le comunicazioni rapide viene utilizzato Telegram in modo da essere sempre disponibile, andrà usato con cautela e parsimonia.\end{itemize}

\item \textbf{textsl{Meetings}}
\begin{itemize} \item Per i meetings si utilizzerà probabilmente discord per la semplicità di fruizione\end{itemize}

\item \textbf{textsl{VCS}}
\begin{itemize} \item Il VCS ovvero Version Control System viene utilizzato per tracciare, mantenere, versionare e gestire tutto il prodotto del progetto, tra cui anche la documentazione.\end{itemize}

\item \textbf{textsl{Google Slides}}
\begin{itemize} \item Per creare le varie presentazioni, ad esempio per il diario di bordo, verrà utilizzato Google Slides\end{itemize}
\end{itemize}

\subsection{Processo di Pianificazione}

\subsection{Ruoli}

Durante il progetto, ogni membro ricoprirà i seguenti ruoli:
\begin{itemize}
\item Responsabile
\item Amministratore
\item Analista
\item Progettista
\item Programmatore
\item Verificatore
\end{itemize}

Per fare ciò, verrà creato un calendario che riporterà i ruoli dei membri nel tempo, in modo che ognuno ricopra almeno ciascun ruolo per un tempo omogeneo.

Inoltre per semplificare visivamente, nel server Discord verranno utilizzati dei tag rappresentanti i vari ruoli. Quindi ogni membro saprà esattamente chi sta svolgendo un qualsiasi ruolo senza dover consultare il calendario.

\begin{center}
    \begin{tabularx}{\textwidth} {
        | >{\raggedright\arraybackslash}>{\hsize=.4\hsize}X |
          >{\raggedright\arraybackslash}>{\hsize=.6\hsize}X |
        }

        \hline
        \textbf{Tag (ruolo)} & \textbf{Colore} \\
        \hline\hline

        %Tag (ruolo)
        Responsabile &
        %Colore
        Rosso \\
        \hline

        %Tag (ruolo)
        Amministratore &
        %Colore
        Arancione \\
        \hline

        %Tag (ruolo)
        Analista &
        %Colore 
        Giallo \\
        \hline

        %Tag (ruolo)
        Progettista &
        %Colore
        Verde \\
        \hline

        %Tag (ruolo)
        Programmatore &
        %Colore
        Blu \\
        \hline

        %Tag (ruolo)
        Verificatore &
        %Colore   
        Violetto \\
        \hline

    \end{tabularx}
\end{center}


(spiegazione dettagliata dei vari ruoli?)

\subsection{Assegnazione dei compiti}

Per l'assegnazione dei compiti, si è deciso di utilizzare le \textbf{textsl{issue}} di GitHub.

\subsubsection{Emettere una issue}

- Per emettere una issue è importante che la issue sia di tipo "azione richiesta"
- La issue non deve contenere nel titolo informazioni relative alla repo in cui si trova
- È bene che la issue abbia delle \textsl{label} in modo da renderla più chiara, facile da riconoscere e trovare 

\subsubsection{Gestione delle issue}

La gestione delle issue è demandata al \textbf{\textsl{responsabile}}, il quale avrà come compito quello di assegnare ad ogni componente del gruppo le issue da risolvere o da completare entro una determinata data.

L'emissione delle issue di bug o problematiche potrà invece avvenire da ogni componente del gruppo. Infatti, chiunque si accorga che qualcosa non va o andrebbe risolta può aggiungere issue in modo che il \textbf{\textsl{responsabile}} possa poi decidere con che modalità procedere.

\subsubsection{Template di issue}

Per veloccizzare e semplificare l'emissione di issue si è deciso di creare dei template, in modo da avere una base già pronta per le issue di uno stesso tipo.

\begin{center}
    \begin{tabularx}{\textwidth} {
        | >{\raggedright\arraybackslash}>{\hsize=.4\hsize}X |
          >{\raggedright\arraybackslash}>{\hsize=.6\hsize}X |
        }

        \hline
        \textbf{Template} & \textbf{Significato} \\
        \hline\hline

        %Template
        Bug report &
        %Significato
       Template per la creazione di issue relativi a bug  \\
        \hline

        %Template
        TODO &
        %Significato
       Template per la creazione di una issue di TODO  \\
        \hline

    \end{tabularx}
\end{center}

GitHub fornisce delle \textsl{label} da utilizzare per raggruppare le issue, ma abbiamo deciso di crearne alcune di personalizzate per rendere ancora più esplicativo il compito di una issue.

\begin{center}
    \begin{tabularx}{\textwidth} {
        | >{\raggedright\arraybackslash}>{\hsize=.4\hsize}X |
          >{\raggedright\arraybackslash}>{\hsize=.6\hsize}X |
        }

        \hline
        \textbf{Label} & \textbf{Significato} \\
        \hline\hline

        %Labe
        Piano di progetto &
        %Significato
       Miglioramenti o aggiunte al piano di progetto  \\
        \hline

        %Label
        Piano di qualifica &
        %Significato
       Miglioramenti o aggiunte al piano di qualifica  \\
        \hline

    \end{tabularx}
\end{center}
